\section{課題}

以下の 8 つの課題を満たす Python プログラムを作成し、実行結果を確認する。

\begin{enumerate}
\item \textbf{動画の読み込み}: 動画(People - 84973.mp4)を読み込み、動画の縦・横のサイズ、
      FPS、総フレーム数を出力する。
\item \textbf{サムネイル画像の作成}: サムネイルとして動画(People - 84973.mp4)中の
      フレーム画像を 1 枚出力し、保存する。ただし、サムネイルは単に 1 フレーム目ではなく、
      中央のフレーム、または自分で良いと思った選択方法で選んだ代表的なフレームとする。
\item \textbf{ファイル名のパース}: 動画のファイル名(People - 84973.mp4)の文字列を処理して、
      タイトル(People)と ID(84973)に分けるプログラムを作成し、その結果を出力する。
      他の動画(Cat - 66004.mp4 など)でも同様のアルゴリズムで分けられるようにする。
\item \textbf{色反転と動画の保存}: 動画(Cat - 66004.mp4)の各フレームの色を反転した動画を
      作成し、保存する。反転後の画素値は「画素値の最大値 $-$ 元の画素値」とする。
      縦・横のサイズは元のサイズからそれぞれ 0.3 倍する。保存するファイル名は
      6324034\_cat\_reverse.mp4 とする。{\tt cv2.bitwise\_not()} は使用しない。
\item \textbf{音声の可視化}: 音声(report\_audio.wav)には時間-周波数領域にメッセージが
      隠されている。パラメータを適宜調節して解析し、可視化したグラフを描画する。
\item \textbf{n-gram}: テキスト(text\_english.txt)に出現する 2-gram の出現頻度を調べ、
      上位 10 個をグラフ化する(6-a)。さらに stop words を設計・除去したうえで
      上位 10 個の 2-gram をグラフ化し、除去しない場合と比較する(6-b)。
\item \textbf{動画の変わり目検出}: 4 つの動画を結合した動画(Detection.mp4)から、
      グレースケール化した隣接フレーム間の画素値差の総和(変化量)を用いて変わり目
      (3 箇所)を自動検出する。横軸フレーム・縦軸変化量のグラフを描画し(7-a)、
      閾値の決め方を 2 種類以上試して変わり目フレームを出力する(7-b)。
      {\tt absdiff()} などの関数は使用しない(numpy は使用可)。
\item \textbf{オブジェクト指向}: 動画のパスを受け取ってインスタンス化する Video クラスを
      作成する。インスタンス変数(パス・タイトル・ID・サムネイル)、課題 2・3・4 の処理の
      メソッド化、インスタンス変数を辞書で返す {\tt get\_info()}、およびフレーム処理関数
      (二値化・グレースケール化・色反転)を引数として受け取る高階関数方式の動画処理
      メソッドを実装し、動作を確認する。
\end{enumerate}
